%/*% pdflatex -file-line-error -interaction nonstopmode %

\documentclass{article}
\usepackage{graphicx}
\usepackage{amsmath}
\usepackage{amsfonts}
\usepackage{bookmark}
\usepackage{exercises}
\usepackage{amsthm}
\usepackage{yhmath}
\newcommand*{\oldneg}{\mathord{\sim}}

\theoremstyle{definition}
\newtheorem{definition}{Definition}

\setlength{\parindent}{.5in}
\setlength{\parskip}{12pt}

\title{1.3 Exercises}
\begin{document}
\maketitle

\section*{1}
The zeroes of a function are all those $x$ such that $f(x)=0$.

\section*{2}
A monomial an expression that is made of one term, that is, a variable, number, or a product of number and one or more variables. A binomial is an expression
 that is the sum two monomials, while a trinomial is an expression that is the sum of three monomials. An example of a monomial is $4x$, an example of a binomial is $4x+3$, and 
an example of a trinomial is $y^2+4x+3$.

\section*{3}
$x^2+6x+5\newline\rightarrow (x+3)(x+2)$

\section*{4}
$x^2-7x+10\newline\rightarrow (x-5)(x-2)$

\section*{5}
$a^2-13a+22\newline\rightarrow (a-11)(a-2)$

\section*{6}
$r^2+15r+56\newline\rightarrow (r+8)(r+7)$

\section*{7}
Cannot be factored.

\section*{8}
$q^2-11q+28\newline\rightarrow (q-7)(q-4)$

\section*{9}
$b^2+3b-40\newline\rightarrow (b+8)(b-5)$

\section*{10}
$x^2-4x-12\newline\rightarrow (x-6)(x+2)$

\section*{11}
$x^2-7x-18\newline\rightarrow (x-9)(x+2)$

\section*{12}
Cannot be factored.

\section*{13}
$x^2+9x-36\newline\rightarrow (x+12)(x-3)$

\section*{14}
$m^2+8m-65\newline\rightarrow (m+13)(m-5)$

\section*{15}
$x^2-36\newline\rightarrow (x+6)(x-6)$

\section*{16}
$b^2-81\newline\rightarrow (b+9)(b-9)$

\section*{17}
$x^2-24x+144\newline\rightarrow (x-12)^2$

\section*{18}
$t^2-16t+64\newline\rightarrow (t-8)^2$

\section*{19}
$x^2+8x+16\newline\rightarrow (x+4)^2$

\section*{20}
$c^2+28c+196\newline\rightarrow (c+14)^2$

\section*{21}
$n^2+14n+49\newline\rightarrow (n+7)^2$

\section*{22}
$s^2-26s+169\newline\rightarrow (s-13)^2$

\section*{23}
$z^2-121\newline\rightarrow (z+11)(z-11)$

\section*{24}
$x^2-8x+12=0\newline\rightarrow (x-6)(x-2)=0\newline\rightarrow x=6\vee x=2$

\section*{25}
$x^2-11x+30=0\newline\rightarrow (x-6)(x-5)=0\newline\rightarrow x=6\vee x=5$

\section*{26}
$x^2+2x-35=0\newline\rightarrow (x+7)(x-5)=0\newline\rightarrow x=-7\vee x=5$

\section*{27}
$a^2-49=0\newline\rightarrow (a+7)(a-7)=0\newline\rightarrow a=-7\vee a=7$

\section*{28}
$b^2-6b+9=0\newline\rightarrow (b-3)^2=0\newline\rightarrow b=3$

\section*{29}
$c^2+5c+4=0\newline\rightarrow (c+4)(c+1)=0\newline\rightarrow c=-4\vee c=-1$

\section*{30}
$n^2-6n=0\newline\rightarrow n(n-6)=0\newline\rightarrow (n-6)(n-0)=0\newline\rightarrow n=6\vee n=0$

\section*{31}
$t^2+10t+25=0\newline\rightarrow (t+5)^2=0\newline\rightarrow t=-5$

\section*{32}
$w^2-16w+48=0\newline\rightarrow (w-12)(w-4)=0\newline\rightarrow w=12\vee w=4$

\section*{33}
$z^2-3z=54\newline\rightarrow z^2-3z-54=0\newline\rightarrow (z-9)(z+6)=0\newline\rightarrow z=9\vee z=-6$

\section*{34}
$r^2+2r=80\newline\rightarrow r^2+2r-80=0\newline\rightarrow (r+10)(r-8)=0\newline\rightarrow r=-10\vee r=8$

\section*{35}
$u^2=-9u\newline\rightarrow u^2+9u=0\newline\rightarrow u(9+u)=0\newline\rightarrow (9+u)(u-0)=0\newline\rightarrow u=-9\vee 9$

\section*{36}
$m^2=7m\newline\rightarrow m^2-7m=0\newline\rightarrow m(m-7)=0\newline\rightarrow (m-7)(m-0)=0\newline\rightarrow m=7\vee m=0$

\section*{37}
$14x-49=x^2\newline\rightarrow x^2-14x+49=0\newline\rightarrow (x-7)^2=0\newline\rightarrow x=7$

\section*{38}
$-3y+28=y^2\newline\rightarrow y^2+3y-28=0\newline\rightarrow (y+7)(y-4)=0\newline\rightarrow y=-7\vee y=4$

\section*{39}
$3-2=1$, not $-1$, meaning that the correct factorization is $(x+2)(x-3)$, and hence, that $x=-2\vee x=3$

\section*{40}
We need to write the function in intercept form in order to find the zeros; we should not solve the function to find values that equal $14$. Instead:\newline
$x^2+7x-8=0\newline\rightarrow (x+8)(x-1)=0\newline\rightarrow x=-8\vee x=1$

\section*{41}
$x^2+2x-63=0\newline\rightarrow (x+9)(x-7)=0\newline\rightarrow x=-9\vee x=7$\newline Hence, the answers is $A$.

\section*{42}
$2(24\times 10)=(24+x)(10+x)\newline\rightarrow 480=x^2+34x+240\newline\rightarrow x^2+34x-240=0\newline\rightarrow (x+40)(x-6)=0\newline\rightarrow 
x=-40\vee x=6$. Becuase we can only use positive values, $x=6$.

\section*{43}
$3(10\times 12)=(12+x)(10+x)\newline\rightarrow 360=x^2+22x+120\newline\rightarrow x^2+22x-240=0\newline\rightarrow (x+30)(x-8)=0\newline\rightarrow 
x=-30\vee x=8$. BEcause we can only use positive values, $x=8$.

\section*{44}
$x^2+6x+8=0\newline\rightarrow (x+4)(x+2)=0\newline\rightarrow x=-4\vee x=-2$

\section*{45}
$x^2-8x+16=0\newline\rightarrow (x-4)(x-4)=0\newline\rightarrow x=4$

\section*{46}
$x^2-4x-32=0\newline\rightarrow (x-8)(x+4)=0\newline\rightarrow x=8\vee x=-4$

\section*{47}
$x^2+7x-30=0\newline\rightarrow (x+10)(x-3)=0\newline\rightarrow x=-10\vee x=3$

\section*{48}
$x^2+11x=0\newline\rightarrow x(x+11)=0\newline\rightarrow (x+11)(x-0)=0\newline\rightarrow x=-11\vee x=0$

\section*{49}
$x^2-8x=0\newline\rightarrow (x-0)(x-8)=0\newline\rightarrow x=0\vee x=8$

\section*{50}
$x^2-64=0\newline\rightarrow (x-8)(x+8)=0\newline\rightarrow x=8\vee x=-8$

\section*{51}
$x^2-25=0\newline\rightarrow (x-5)(x+5)=0\newline\rightarrow x=5\vee x=-5$

\section*{52}
$x^2-12x-45=0\newline\rightarrow (x-15)(x+3)=0\newline\rightarrow x=15\vee x=-3$

\section*{53}
$x^2+19x+84=0\newline\rightarrow (x+12)(x+7)=0\newline\rightarrow x=-12\vee x=-7$

\section*{54}
$x^2+22x+121=0\newline\rightarrow (x+11)(x+11)=0\newline\rightarrow x=-11$

\section*{55}
$x^2+2x+1=0\newline\rightarrow (x+1)(x+1)=0\newline\rightarrow x=-1$

\section*{56}
$x^2+6x-55=0\newline\rightarrow (x+11)(x-5)=0\newline\rightarrow x=-11\vee x=5$ Hence, the answers is $B$.

\section*{57}
$x^2-19x+88$

\section*{58}
For the expression $x^2+bx+7$ to be factored, $b=m+n$ where $mn=7$. Hence, the only two such cases are for $m=7, n=1$, and $m=-7, n=-1$. Thus, 
the formula $x^2+bx+7$ can be factored iff $b=8\vee b=-8$.

\section*{59}
$36=(x+5)x\newline\rightarrow x^2+5x-36=0\newline\rightarrow (x+9)(x-4)=0\newline\rightarrow x=-9\vee x=4$, and thus, because only positive values are 
acceptable, $x=4$.

\section*{60}
$84=(x+7)(x+2)\newline\rightarrow x^2+9x+14=84\newline\rightarrow x^2+9x-70=0\newline\rightarrow (x+14)(x-5)=0\newline\rightarrow x=5\vee x=-14$, and thus, 
because only positive values are acceptable, $x=5$.

\section*{61}
$42=\frac{1}{2}(x+3)(2x+8)\newline\rightarrow \frac{1}{2}(2x^2+14x+24)=42\newline\rightarrow x^2+7x+12=42\newline\rightarrow
x^2+7x-30=0\newline\rightarrow (x+10)(x-3)=0\newline\rightarrow x=-10\vee x=3$, and thus, because negative values are ruled out, $x=3$. 

\section*{62}
$32=\frac{1}{2}(x+6+x+2)x\newline\rightarrow \frac{1}{2}(2x^2+8x)=32\newline\rightarrow x^2+4x=32\newline\rightarrow x^2+4x-32\newline\rightarrow 
(x+8)(x-4)=0\newline\rightarrow x=-8\vee x=4$, and thus, because negative values are ruled out, $x=4$.

\section*{63}
We will choose the values $0$ and $20$, both $10$ away from $10$ on a number line:
$x^2-20x$

\section*{64}
(a):\newline For $x^2+16=(x+m)(x+n)$ to be true, $x^2+16=x^2+x(m+n)+mn$ would be true. Now, this means that $m+n=0$ and $mn=16$ because the coefficients 
of the corresponding terms of polynomials are equal if two polynomials are equal.\newline 
(b):Because these can be only true if $m$ and $n$ have opposite signs, which 
cannot be the case because $mn=16$, these two equations cannot both be true. Hence, this equation cannot be factored by any $m,n$. 

\section*{65}
$3(100\times 50)=(100+x)(50+x)\newline\rightarrow 15000=x^2+150x+5000\newline\rightarrow x^2+150x-10000=0\newline\rightarrow (x+200)(x-50)=0\newline
\rightarrow x=-200\vee x=50$, and thus the dimensions of the new park are $150$ by $100$ ft. wide. 

\section*{66}
$2(35\times 18)=(35+x)(18+x)\newline\rightarrow 1260=x^2+53x+630\newline\rightarrow x^2+53x-630=0\newline\rightarrow (x+63)(x-10)=0\newline\rightarrow 
x=-63\vee x=10$, and thus, the new dimensions of the enclosure are $45$ feet long by $28$ feet wide.

\section*{67}
(a): $A=20\times 30=600 ft.^2$\newline
(b): $600+464=(30+x)(20+x)$\newline
(c): $1064=x^2+50x+600\newline\rightarrow x^2+50x-464=0\newline\rightarrow (x+58)(x-8)=0\newline\rightarrow x=-58\vee x=8$, and thus, the length and the 
width of the of the patio should be expanded by $8$ ft. 

\section*{68}
(a): $x^2+5x+6$\newline
(b): $(x+3)(x+2)$ The model an $x$ grouped with $2$ ones, and an $x$ grouped with $3$ ones, which is similar to $x+2$ and $x+3$, the terms of the 
factorization.

\section*{69}
$3(18\times 15)=(18+x)(15+x)\newline\rightarrow 810=x^2+33x+270\newline\rightarrow x^2+33x-540=0\newline\rightarrow (x+45)(x-12)=0\newline\rightarrow 
x=-45\vee x=12$, and thus, the length of the new rope needed is $12+12=24$ ft.

\section*{70}
$\frac{21\times 20}{2}=(21-x)(20-x)\newline\rightarrow 210=x^2-41x+420\newline\rightarrow x^2-41x+210=0\newline\rightarrow (x-35)(x-6)=0\newline\rightarrow 
x=35\vee x=6$, because negative measurements are ruled out, $x=6$. Thus, the new deck is $15$ ft. by $14$ ft.

\section*{71}
$200=(10+x)(10+x)\newline\rightarrow x^2+20x-100=0\newline\rightarrow$ Because $x^2+20x-100$ cannot be factored, 
the problem cannot be solved by factoring. 

\section*{72}
This problem can be modeled with an equation that increases the width and length of the old lot and grocery store both by $x$, and then subtracting the 
area of the grocery store from the total area obtained by solving the quadratic equation:\newline
$(240+x)(375+x)=130500\newline\rightarrow x^2+615x+90000=130500\newline\rightarrow x^2+615-40500=0\newline\rightarrow (x+675)(x-60)=0\newline\rightarrow 
x=-675\vee x=60$, and thus, because negative values are ruled out, $x=60$.
\end{document}
 
